\section*{الفصل \ref{ch:b3:advanced-nmr} --- الرنين المغناطيسي النووي بعمق: النبضات والاسترخاء والأطياف الثنائية البعد}

\begin{solution}{exo:b3:advanced-nmr:1}
$\nu_0 = (\gamma/2\pi)B_0$: \ce{^1H} \qty{600.4}{MHz}، و\ce{^{13}C} \qty{151.0}{MHz}، و\ce{^{19}F}
\qty{565.1}{MHz}.
\end{solution}

\begin{solution}{exo:b3:advanced-nmr:2}
$h\nu_0/2kT = 6.626\times10^{-34} \times 400.2\times10^6/(2 \times 1.381\times10^{-23} \times 300)
= \num{3.2e-5}$. وعند \qty{4}{K} تكون أكبر بخمس وسبعين مرة، \num{2.4e-3} (والتقريب
ما زال صالحًا).
\end{solution}

\begin{solution}{exo:b3:advanced-nmr:3}
$\theta = \gamma B_1t_p = \pi/2$ في \qty{10}{\micro s}: $\gamma B_1/2\pi = 1/(4 \times
\qty{10}{\micro s}) = \qty{25}{kHz}$. وتدوم النبضة $180^\circ$ مدة \qty{20}{\micro s}؛ والنبضة \qty{5}{\micro s}
تعطي $45^\circ$.
\end{solution}

\begin{solution}{exo:b3:advanced-nmr:4}
$T_2 = 1/(\pi \times 3.2) = \qty{0.10}{s}$.
\end{solution}

\begin{solution}{exo:b3:advanced-nmr:5}
$(10.71/42.58)^3 = 0.0159$، مضروبًا في 0.011: \num{1.75e-4}، أي نحو 1/5700 من إشارة
البروتون. وبما أن نسبة الإشارة إلى الضجيج $\propto\sqrt N$، فالنسبة نفسها تقتضي عددًا من المسوح أكبر بنحو $5700^2 \approx \num{3e7}$
مرة --- وهذا مستحيل، ولهذا تستعمل أطياف \babelsublr{${}^{13}$C} عيّنات أكثر تركيزًا،
ونزع الاقتران ونقل الاستقطاب.
\end{solution}

\begin{solution}{exo:b3:advanced-nmr:6}
$T_1 = 1.39/\ln2 = \qty{2.0}{s}$؛ وفترة التكرار $\ge 5T_1 = \qty{10}{s}$.
\end{solution}

\begin{solution}{exo:b3:advanced-nmr:7}
\omterm{def:b3:advanced-nmr:experiments}{COSY}: \omterm{def:b3:advanced-nmr:two-d}{قمة متقاطعة} واحدة، بين رباعي \ce{CH2} وثلاثي \ce{CH3} لمجموعة
الإيثيل؛ وأحادي \ce{CH3CO} لا يقترن بشيء. \omterm{def:b3:advanced-nmr:experiments}{وHMBC} انطلاقًا من الميثيل الأحادي:
كربون الكربونيل (رابطتان) وكربون \ce{CH2} (ثلاث روابط).
\end{solution}

\begin{solution}{exo:b3:advanced-nmr:8}
$k_c = \pi \times 50/\sqrt2 = \qty{111}{s^{-1}}$. $\Delta G^\ddagger = 8.314 \times 330 \times
\ln(6.88\times10^{12}/111) = \qty{68}{kJ/mol}$.
\end{solution}

\begin{solution}{exo:b3:advanced-nmr:9}
NOE $\propto r^{-6}$: $r_b = 0.30 \times 4^{-1/6} = \qty{0.24}{nm}$.
\end{solution}

\begin{solution}{exo:b3:advanced-nmr:10}
في \omterm{def:b3:advanced-nmr:magnetisation}{المعلم الدوّار} عند الرنين لا يبقى إلا $\mathbf B_1 = B_1\mathbf e_{x'}$:
$\dot M_{x'} = 0$، $\dot M_{y'} = \gamma B_1M_z$، $\dot M_z = -\gamma B_1M_{y'}$ (مركبات
$\gamma\mathbf M\times\mathbf B_1$). فتبقى $M_{x'}$ ثابتة ويدور $(M_{y'}, M_z)$ بالسرعة الزاوية
$\gamma B_1$: فبعد $t_p$، بالزاوية $\gamma B_1t_p$ حول $x'$.
\end{solution}

\begin{solution}{exo:b3:advanced-nmr:11}
السبين الواقع في مجال أقوى قليلًا يكتسب طورًا بمعدل إضافي ثابت خلال
$\tau$؛ والنبضة $180^\circ$ تعكس طوره، فيخسر المقدار نفسه بالضبط
خلال $\tau$ الثانية: فتعود كل هذه السبينات إلى الاصطفاف عند $2\tau$. أما التآثرات
العشوائية بين السبينات فتتقلب مع الزمن ولا تتكرر في الفترة الثانية: فلا
يُلغى ما تسببه من فقدان للتطاور. وتتناقص سعة الصدى بالزمن $T_2$ الحقيقي.
\end{solution}

\begin{solution}{exo:b3:advanced-nmr:12}
إيثانوات الإيثيل، \ce{CH3COOCH2CH3}: رباعي \ce{OCH2} (4.12)، وأحادي \ce{CH3CO} (2.05)،
وثلاثي \ce{CH3}؛ وبروتونات الرباعي ترى كربون الكربونيل على بعد ثلاث روابط،
عبر أكسجين الإستر. أما بروبانوات الميثيل، \ce{CH3CH2COOCH3}، فتُظهر
رباعيها قرب 2.3 ppm (على ذرة كربون قرب 27 ppm في \omterm{def:b3:advanced-nmr:experiments}{HSQC}، لا 60) وأحاديًا
\ce{OCH3} قرب 3.7 ppm تعبر قمته في \omterm{def:b3:advanced-nmr:experiments}{HMBC} نحو الكربونيل ذرة الأكسجين.
\end{solution}

\begin{solution}{pb:b3:advanced-nmr:1}
\textbf{1.} $(2 \times 6 + 2 - 12)/2 = 1$: رابطة C=O واحدة.
\textbf{2.} 173.6 ppm: كربونيل الإستر؛ 60.1 ppm: \ce{OCH2}.
\textbf{3.} خمس؛ DEPT-135: \ce{CH2} سالبة (60.1، 36.3، 18.6)، و\ce{CH3} موجبة (14.3،
13.7)، والكربونيل غائب.
\textbf{4.} \qty{400.2}{MHz} و\qty{100.7}{MHz}.
\textbf{5.} $3J = \qty{21.6}{Hz} = \qty{0.054}{ppm}$.
\textbf{6.} التعدديات تبيّن أي المجموعات متجاورة داخل كل
شظية، لكنها لا تبيّن كيف تتصل الشظايا عبر الأكسجين والكربونيل.
\textbf{7.} 4.13--1.25؛ 2.28--1.66؛ 1.66--0.95 (والقمم المتناظرة معها).
\textbf{8.} \ce{OCH2CH3} (4.13، 1.25) و\ce{CH2CH2CH3} (2.28، 1.66، 0.95).
\textbf{9.} يفصل بين بروتونات \ce{OCH2} وبروتونات \ce{CH2C=O} الأكسجين 
وكربون الكربونيل: خمس روابط، ولا اقتران مُميَّز.
\textbf{10.} 1.66 ppm، بين \ce{CH2} و\ce{CH3}: مقترن بعدد 2 + 3 = 5 بروتونات بقيم
$J$ متقاربة، فهو سداسي (متعدد خطوط).
\textbf{11.} ميثيلان مقترنان، أي مجموعتا \ce{CH3} على ذرتي كربون متجاورتين (وحدة \ce{CH3-CH3}): مستحيل هنا.
\textbf{12.} \ce{-O-CH2-CH3} و\ce{-CH2-CH2-CH3}، يضاف إليهما \ce{C=O}.
\textbf{13.} 4.13/60.1؛ 2.28/36.3؛ 1.66/18.6؛ 1.25/14.3؛ 0.95/13.7.
\textbf{14.} 14.3 (\ce{CH3} الإيثيل) و13.7 ppm (\ce{CH3} البوتانويل)، ولا يفصل بينهما إلا \qty{0.6}{ppm}.
\textbf{15.} 4.13 ppm نحو 173.6 ppm: H--C--O--C، ثلاث روابط.
\textbf{16.} 2.28 ppm (رابطتان) و1.66 ppm (ثلاث روابط) نحو 173.6 ppm.
\textbf{17.} \ce{CH3CH2CH2C(=O)OCH2CH3}، بوتانوات الإيثيل. أما بروبانوات البروبيل فتضع
ثلاثيًا \ce{OCH2} (لا رباعيًا) قرب 4.0 ppm ورباعيًا قرب 2.3 ppm.
\textbf{18.} الاقترانات عبر أربع روابط عادة دون 1 Hz؛ وفترة انتظار \omterm{def:b3:advanced-nmr:experiments}{HMBC}، المضبوطة على
5--10 Hz، لا تنتقيها.
\textbf{19.} تسترخي ذرات الكربون بسرعات مختلفة وتتلقى تعزيزات أوفرهاوزر
نووية مختلفة من نزع الاقتران.
\textbf{20.} $1 - \eu^{-5/20} = 0.22$: الكربونيل ممثَّل تمثيلًا ناقصًا بعامل 4.
\textbf{21.} $1 - \eu^{-5} = 0.993$.
\textbf{22.} مفعول NOE الناتج عن نزع اقتران البروتونات، المختلف من ذرة كربون إلى أخرى؛ ويُزال
بنزع الاقتران المبوَّب العكسي (جهاز نزع الاقتران لا يعمل إلا أثناء الاقتناء).
\textbf{23.} $256 \times \qty{100}{s} = \qty{25600}{s}$، أي نحو 7 ساعات.
\textbf{24.} \textbf{$5T_1$ لكربون الكربونيل: \qty{100}{s} بين المسوح.}
\end{solution}
